The center must be at least one unit from every cell edge. Its favorable region is a $3$ by $1$ rectangle, giving $\boxed{\dfrac{3}{15}=\dfrac{1}{5}}$.

For the first group,
\[
\sigma=\sqrt{\dfrac{295}{15}-\left(\dfrac{61}{15}\right)^2}
=\boxed{1.76887\ldots}.
\]
For the second group, $\sum x=12(3.5)=42$ and $\sum x^2=12(2^2+3.5^2)=195$. The combined totals are $n=27,\sum x=103,\sum x^2=490$, so
\[
\boxed{\bar x=\dfrac{103}{27}=3.81481\ldots},\qquad
\boxed{\sigma=\sqrt{\dfrac{490}{27}-\left(\dfrac{103}{27}\right)^2}=1.89614\ldots}.
\]
A disc covers a corner exactly when its center lies within one unit of that corner. In a fundamental cell these regions are four disjoint quarter-discs, of total area $\pi$. Thus the probability is $\dfrac{\pi}{15}$ and the experimental estimate is $\boxed{\pi\approx15(0.2216)=3.324}$.
\begin{asy}
texpreamble("\usepackage{amsmath,amssymb}");
size(240);
defaultpen(fontsize(9pt));
fill(box((1,1),(4,2)),lightgreen);draw(box((0,0),(5,3)));draw(box((1,1),(4,2)),blue);draw(arc((0,0),1,0,90),red);draw(arc((5,0),1,90,180),red);draw(arc((5,3),1,180,270),red);draw(arc((0,3),1,270,360),red);label("$5$",(2.5,0),S);label("$3$",(0,1.5),W);label("Prize centers",(2.5,1.5));
\end{asy}
